\chapter{GIỚI THIỆU ĐỀ TÀI}
\label{chap:gioithieu}

\section{Động cơ thực hiện đề tài}
Suy luận mạng nơ-ron tại thiết bị biên cho phép xử lý dữ liệu gần nơi phát sinh. Tuy nhiên, các hệ thống nhúng thường bị giới hạn về năng lực tính toán và dung lượng bộ nhớ. Với mạng nơ-ron tích chập (Convolutional Neural Network, CNN), số lượng phép nhân--cộng và các lần truy cập tensor có thể làm thời gian suy luận trên một CPU nhỏ trở nên lớn. Hiệu quả năng lượng là một động cơ chung của tính toán tại biên; trong đề tài này, các đại lượng được đánh giá là tính đúng đắn, số chu kỳ, tài nguyên FPGA và timing.

Phần cứng chuyên dụng có thể giảm chi phí thực hiện một số phép toán CNN, nhưng cần thêm tài nguyên lưu trữ, datapath và logic điều khiển. Khi CPU phải cấp toán hạng và thu kết quả qua các thanh ghi ánh xạ bộ nhớ, chi phí truyền dữ liệu có thể hạn chế lợi ích của datapath. Vì vậy, lựa chọn phần việc đưa xuống phần cứng phải gắn với cách firmware tổ chức dữ liệu và gọi bộ tăng tốc.

Từ vấn đề đó, đề tài hiện thực MyTinyMobileNetSoC, một hệ thống dựa trên PicoRV32 để suy luận TinyMobileNet trên dữ liệu MNIST. Workload có các lớp tích chập chuẩn, depthwise và pointwise, cho phép khảo sát những yêu cầu khác nhau về tính toán và tái sử dụng dữ liệu trong cùng một đường suy luận. Trọng tâm là xây dựng hệ thống hoạt động tích hợp, xác định nút thắt bằng số đo và cải thiện hiệu năng qua đồng thiết kế phần cứng--phần mềm (HW/SW co-design).

\section{Bài toán đặt ra}
Tăng tốc CNN không chỉ là tăng tốc số học. Trong MyTinyMobileNetSoC, CPU đọc tensor từ RAM, ghi toán hạng qua MMIO, điều khiển thực thi và đọc kết quả về. CPU còn điều phối thứ tự lớp, cộng bias, xử lý offset và giới hạn giá trị đầu ra. Các bộ tăng tốc không tự truy cập RAM và không có giao diện DMA.

Do đó, hiệu năng phụ thuộc đồng thời vào điều khiển CPU, số giao dịch MMIO, khả năng tái sử dụng dữ liệu, lịch vòng lặp, tái lượng tử hóa và chi phí mỗi lần gọi bộ tăng tốc. Một khối tính toán nhanh vẫn có thể dành phần lớn thời gian chờ CPU cấp dữ liệu. Ngược lại, giữ lại đầu vào hoặc trọng số qua nhiều phép tính có thể giảm chu kỳ mà không cần nhân rộng datapath.

Bài toán của luận văn là phân chia và phối hợp công việc giữa PicoRV32 với các bộ tăng tốc để giảm số chu kỳ suy luận, đồng thời kiểm tra kết quả và đánh giá chi phí hiện thực FPGA. Việc tối ưu được thực hiện trên cả firmware lẫn phần cứng, với đường xử lý trên CPU được duy trì cho các trường hợp chưa được tăng tốc.

\section{Mục tiêu của đề tài}
Các mục tiêu được xác định theo những kết quả có thể kiểm tra:
\begin{enumerate}
    \item Tích hợp suy luận TinyMobileNet trên PicoRV32, từ dữ liệu ảnh MNIST và tham số mô hình đến điểm số của 10 lớp và nhãn dự đoán.
    \item Hiện thực và tích hợp hỗ trợ phần cứng chuyên dụng cho tích chập chuẩn (Conv), tích chập theo chiều sâu (Depthwise), tích chập theo điểm (Pointwise) và tái lượng tử hóa (Quant).
    \item Duy trì firmware điều phối cùng đường xử lý thay thế trên CPU (fallback), bảo đảm các lớp hoặc cấu hình ngoài phạm vi bộ tăng tốc vẫn có đường thực thi.
    \item Giảm số chu kỳ của đường suy luận tích hợp bằng tối ưu truyền dữ liệu, lịch vòng lặp và tái sử dụng toán hạng; ghi nhận kết quả qua sáu mức tối ưu.
    \item Đánh giá tính đúng đắn trong phạm vi các phép thử đã thực hiện; so sánh điểm số, nhãn dự đoán và số chu kỳ giữa các phiên bản, đồng thời phân biệt baseline tăng tốc với đối chứng CPU-only.
    \item Đánh giá tài nguyên sử dụng và timing sau định tuyến (post-route) của thiết kế cuối trên FPGA mục tiêu ZCU106.
\end{enumerate}

\section{Phạm vi thực hiện}
\subsection{Workload và biểu diễn số}
Đối tượng suy luận là TinyMobileNet tùy biến, lấy cảm hứng từ cách tổ chức depthwise--pointwise của MobileNet; đây không phải MobileNetV2. Mạng phân loại ảnh MNIST thang xám kích thước $28\times28$ thành 10 lớp. Chuỗi lớp gồm Conv0, ba cặp depthwise--pointwise, trung bình toàn cục (GAP) và lớp kết nối đầy đủ (FC). Firmware chọn nhãn từ điểm số đầu ra bằng argmax.

Kích hoạt và trọng số dùng số nguyên có dấu INT8; bias và giá trị tích lũy dùng INT32. Tái lượng tử hóa sử dụng số học điểm cố định để đưa kết quả về miền đầu ra. Phạm vi kiểm chứng phải xét ngữ nghĩa làm tròn, offset và giới hạn giá trị của đường hiện thực; không mặc nhiên giả định tương đương từng bit với TFLite.

\subsection{Nền tảng phần cứng và phần mềm}
Hệ thống dùng lõi PicoRV32, firmware bare-metal và RAM dùng chung cho mã lệnh, dữ liệu với dung lượng 256~KiB. CPU trao đổi với các bộ tăng tốc qua MMIO và trực tiếp quản lý di chuyển dữ liệu. SoC còn có UART, thanh ghi LED và bộ nhân độc lập; bộ nhân độc lập không được đường suy luận TinyMobileNet sử dụng.

Phần cứng hỗ trợ tích lũy cho Conv0, các nhánh Depthwise được hỗ trợ và các lớp Pointwise. Ở phiên bản cuối, DW2 tái sử dụng bộ tăng tốc Depthwise, còn DW3, GAP và FC vẫn tích lũy trên CPU; các lớp này sử dụng phần cứng tái lượng tử hóa. Cách phân công cụ thể được trình bày trong Chương~\ref{chap:kientruc}.

\subsection{Đánh giá FPGA và giới hạn triển khai}
Nền tảng đánh giá FPGA là ZCU106, với linh kiện mục tiêu XCZU7EV. Công việc tổng hợp, bố trí, định tuyến và đánh giá tài nguyên/timing của thiết kế cuối đã hoàn thành trong flow đánh giá hiện có. Điểm ràng buộc chu kỳ $8{,}875$~ns là điểm đạt timing có tần số cao nhất trong các điểm đã thử và có kết luận; kết quả này không xác định Fmax thực sự.

Flow đánh giá SoC chưa khởi tạo RAM bằng firmware và chưa tạo bitstream suy luận TinyMobileNet. Vì vậy, suy luận TinyMobileNet trên bo mạch vật lý vẫn chưa được xác nhận. Phạm vi kết quả không bao gồm đo năng lượng, hiện thực ASIC hoặc chứng minh vận hành thời gian thực trong ứng dụng IoT thực tế.

\section{Phương pháp thực hiện}
Đề tài bắt đầu từ đường suy luận tích hợp đã sử dụng các bộ tăng tốc, làm baseline cho quá trình tối ưu. CPU giữ vai trò điều phối; mỗi thay đổi được xem xét theo khả năng giảm điều khiển, giảm truyền dữ liệu hoặc tăng tái sử dụng, cùng với yêu cầu kiểm tra kết quả. Sáu mức tối ưu lần lượt là:
\begin{enumerate}
    \item \textbf{Level 1:} tối ưu vòng lặp tải trọng số Pointwise trong firmware để giảm chi phí tính địa chỉ và điều khiển vòng lặp.
    \item \textbf{Level 2:} đóng gói bốn trọng số INT8 trong một lần ghi MMIO 32 bit của Pointwise.
    \item \textbf{Level 3:} gom điều khiển Pointwise theo nhóm bốn kênh đầu ra tại một vị trí không gian, dùng chung datapath và cơ chế chờ kết quả trên bus.
    \item \textbf{Level 4:} tái sử dụng phần cứng Depthwise cho DW2 bằng cách chọn các kết quả phù hợp với stride~2.
    \item \textbf{Level 5:} bổ sung cache trọng số cục bộ cho PW1 để giữ trọng số qua các vị trí không gian trong một lần thực thi lớp.
    \item \textbf{Level 6:} đổi lịch Conv0 sang duyệt tile trước (tile-first), giữ đầu vào của tile để dùng lại qua các kênh đầu ra.
\end{enumerate}
Các mức này tạo thành quá trình tối ưu ở cấp hệ thống, thay vì chỉ cải thiện riêng số phép tính mỗi chu kỳ. Kiến trúc và đường suy luận được mô tả trong Chương~\ref{chap:kientruc}; thay đổi, số đo và đánh đổi của từng mức được phân tích trong Chương~\ref{chap:toiuu}.

Việc đánh giá được tách theo cấp bằng chứng. Độ chính xác mô hình và native C trên máy chủ được đánh giá trên 10\,000 ảnh MNIST test. Mô phỏng RTL/SoC tích hợp kiểm tra điểm số, nhãn và các trường hợp tensor hoặc giao dịch đã thử; kết quả này không phải độ chính xác RTL trên toàn bộ tập test. Các mốc chu kỳ Level~1--6 là số đo mô phỏng lịch sử được tổng hợp ở Chương~\ref{chap:toiuu}, không phải phép đo trên bo mạch.

Khoảng đo tổng chu kỳ có chứa UART giữa các lớp và chi phí điều phối, nhưng chưa bao gồm argmax và báo cáo kết quả cuối. Báo cáo FPGA post-route được dùng để đánh giá tài nguyên và timing theo ràng buộc của từng run, tách biệt với kiểm chứng suy luận vật lý. Việc thiếu bộ log và danh tính build lịch sử đầy đủ cũng giới hạn khả năng tái lập đồng thời mọi mức tối ưu.

\section{Các đóng góp của đề tài}
Cần phân biệt thành phần kế thừa với phần công việc của sinh viên. PicoRV32 là lõi mã nguồn mở của bên thứ ba~\cite{YosysHQ2024}. Các khối Conv và Quant được kế thừa từ mã dự án trước: hai tệp \texttt{my\_conv\_accel.v} và \texttt{my\_quant\_accel.v} trong target hiện tại trùng nội dung với các tệp tương ứng trong \texttt{MyLeNet5SoC}; luận văn không nhận quyền tác giả ban đầu đối với các khối này. TensorFlow/Keras, TFLite, dữ liệu MNIST và các công cụ biên dịch, mô phỏng, FPGA là thành phần bên ngoài được sử dụng.

Phần công việc của sinh viên trong đề tài tập trung vào tích hợp MyTinyMobileNetSoC và đường firmware suy luận TinyMobileNet, tích hợp DW/PW cùng giao diện thanh ghi và MMIO, thực hiện tối ưu Level~1--6 (bao gồm PW batch/cache, tái sử dụng phần cứng cho DW2 và lịch Conv0 tile-first), tổ chức script/quy trình đánh giá trong target hiện tại và chuẩn bị đánh giá, triển khai FPGA. Phạm vi này không bao hàm quyền tác giả đối với lõi CPU, các khối kế thừa hoặc framework/công cụ bên ngoài. Các đóng góp nằm trong phạm vi hiện thực và đánh giá hệ thống:
\begin{itemize}
    \item \textbf{Tích hợp hệ thống:} xây dựng đường suy luận TinyMobileNet INT8 trên PicoRV32 với các bộ tăng tốc Conv, Quant, Depthwise và Pointwise giao tiếp qua MMIO do CPU điều khiển.
    \item \textbf{Đồng thiết kế phần cứng--phần mềm:} xác định phân công toán tử, duy trì fallback trên CPU và tổ chức tường minh các bước tái lượng tử hóa, hậu xử lý cùng di chuyển dữ liệu.
    \item \textbf{Tối ưu có số đo:} thực hiện sáu mức tối ưu, giảm baseline đã tăng tốc từ 62\,071\,887 xuống 26\,523\,121 chu kỳ, tương ứng giảm 57,27\% và tăng tốc khoảng $2{,}34\times$ trong khoảng đo đã nêu.
    \item \textbf{Đối chiếu baseline:} so sánh Level~6 với đối chứng CPU-only dùng các kernel C tổng quát trên RV32, có 259\,218\,212 chu kỳ, cho tỷ số khoảng $9{,}77\times$. Đây là lợi ích kết hợp của lịch phần mềm và offload phần cứng; đối chứng chưa được tối ưu chuyên sâu cho CPU.
    \item \textbf{Đánh giá FPGA:} tổng hợp kết quả tài nguyên và timing post-route của thiết kế cuối cho ZCU106, đồng thời chỉ rõ giới hạn của flow và phần kiểm chứng trên bo mạch còn chờ thực hiện.
\end{itemize}
Các kết quả giữ nguyên điểm số và nhãn trên ảnh đối chiếu lịch sử, cùng các kiểm tra tập trung được ghi nhận, hỗ trợ kết luận trong phạm vi đã thử. Chúng không cấu thành chứng minh tương đương mọi đầu vào giữa các phiên bản hoặc giữa RTL, native C và TFLite.

\section{Bố cục luận văn}
Sau chương giới thiệu, nội dung luận văn được tổ chức như sau:
\begin{itemize}
    \item \textbf{Chương 2} trình bày cơ sở về các phép toán CNN, lượng tử hóa, RISC-V/PicoRV32, giao tiếp MMIO, tăng tốc trên FPGA và các thuật ngữ đánh giá hiệu năng.
    \item \textbf{Chương 3} khảo sát các kiến trúc tăng tốc liên quan và đối chiếu với cách tiếp cận của đề tài.
    \item \textbf{Chương 4} mô tả kiến trúc MyTinyMobileNetSoC đã hiện thực, các bộ tăng tốc và đường suy luận phối hợp phần cứng--phần mềm.
    \item \textbf{Chương 5} trình bày quá trình tối ưu, đánh giá tính đúng đắn và hiệu năng theo chu kỳ, cùng kết quả tài nguyên và timing FPGA.
\end{itemize}
